\section*{Chapter \ref{ch:rc:multi-curve-and-collateral-discounting} --- Multi-Curve and Collateral Discounting}

\begin{solution}{exo:rc:multi-curve-and-collateral-discounting:1}
$2.70-2.55 = 15$ basis points
(\cref{prop:rc:multi-curve-and-collateral-discounting:basisrates}), added to the
overnight leg: the party receiving six-month Euribor pays €STR compounded plus
15 basis points.
\end{solution}

\begin{solution}{exo:rc:multi-curve-and-collateral-discounting:2}
The single-curve argument needs a six-month loan to cost the same as two
three-month loans or 180 overnight loans; a persistent spread of fifty
basis points between unsecured and secured term lending made each tenor a
different price. A floating leg on one index was no longer worth par on a curve
built from another, and discounting at a bank's term rate was no longer the
cost of a collateralised position.
\end{solution}

\begin{solution}{exo:rc:multi-curve-and-collateral-discounting:3}
The €STR curve discounts; the six-month Euribor \omterm{def:rc:multi-curve-and-collateral-discounting:curves}{projection curve} projects. With
€STR plus 10 basis points the \omterm{def:rc:multi-curve-and-collateral-discounting:csa}{collateral rate} is higher, so the \omterm{def:rc:multi-curve-and-collateral-discounting:curves}{discount curve}
is the €STR curve shifted by 10 basis points (discount factors times
$e^{-0.001t}$); the \omterm{def:rc:multi-curve-and-collateral-discounting:curves}{projection curve}, recalibrated on it, moves slightly, and
an off-market trade changes value.
\end{solution}

\begin{solution}{exo:rc:multi-curve-and-collateral-discounting:4}
The effective spread is $\max(0,-b(t)) = (15-2t)$ basis points for $t<7.5$
and zero after: $\int_0^{7.5}(15-2t)\,dt = 112.5-56.25 = 56.25$ basis-point
years. The ratio of discount factors is $e^{-0.005625}=0.99439$, and
$68\,598\,292\times(1-0.99439) = \text{USD}~384\,782$.
\end{solution}

\begin{solution}{exo:rc:multi-curve-and-collateral-discounting:5}
The single curve discounts at Euribor rates, about 15 basis points above the
€STR rates at which the collateralised swap is really funded, so it
undervalues the receiver's future net receipts, by EUR~60\,300. Both
frameworks reprice the par ten-year swap by construction, so they agree at the
par rate and differ only for off-market trades, by an amount that grows with
the distance from par.
\end{solution}

\begin{solution}{exo:rc:multi-curve-and-collateral-discounting:6}
A payer of 1.50\% when the par rate is far lower has a liability: its value
is $-\text{EUR}~16.70$ million, and lower discount rates make the liability
larger, a loss of EUR~73\,738, the mirror of the ten-year bar. The clearing
house compensates it: the client receives EUR~73\,738, paid out of the
receivers' gains.
\end{solution}

\begin{solution}{exo:rc:multi-curve-and-collateral-discounting:7}
11.76 basis points, against a ten-year par basis of 15. The par basis is an
annuity-weighted average of the forward bases over ten years; the forward basis
starts at 21.6 and declines, so the ten-year average exceeds the forward at
ten years. Hedging a forward-starting basis position with par basis swaps
needs the whole term structure.
\end{solution}

\begin{solution}{exo:rc:multi-curve-and-collateral-discounting:8}
The poster can switch collateral at any time over the trade's life, so the
option is on the whole path of the basis, not today's value: even with
dollars cheapest today, future dates where euros are better lower the value of
a receivable (the effective rate is the maximum at each date), and volatility
of the basis adds time value. Valuing on the dollar curve ignores both.
\end{solution}

\begin{solution}{pb:rc:multi-curve-and-collateral-discounting:1}
\textbf{1.} An asset: it receives 1.50\% while the par rate is 0.08\%, so its
future net receipts are large.
\textbf{2.} EUR~145\,577\,447.
\textbf{3.} The fixed payments and the current Euribor period are known; later
Euribor fixings are projected from the \omterm{def:rc:multi-curve-and-collateral-discounting:curves}{projection curve}.
\textbf{4.} It has received variation margin equal to the swap's value and
pays \omterm{def:rc:multi-curve-and-collateral-discounting:pai}{price alignment interest} on it at EONIA before the switch, at €STR (8.5
basis points less) after; with negative rates the ``payment'' is a receipt.
\textbf{5.} The fund's position is financed by the variation margin at the PAI
rate, so its value accrues at that rate
(\cref{prop:rc:multi-curve-and-collateral-discounting:csa}).
\textbf{6.} Each discount factor is multiplied by $e^{0.00085\,t}$: 1.0171 at
twenty years.
\textbf{7.} EUR~146\,774\,873.
\textbf{8.} $+\text{EUR}~1\,197\,426$: the fund gains, as a lower discount rate
raises the value of net receipts.
\textbf{9.} The fund pays EUR~1\,197\,426 in cash to the clearing house, which
passes it to the members whose positions lost.
\textbf{10.} The market quotes of Euribor swaps did not move over the weekend;
the method isolates the pure discounting effect, is simple to replicate by each
member, and matches the working group's view that forwards do not depend on
the CSA.
\textbf{11.} EUR~141\,575 per basis point.
\textbf{12.} It was EONIA discounting risk and is now €STR discounting risk.
EONIA was €STR plus a fixed 8.5 basis points, so the two risks were the same
and nothing needed hedging; federal funds and SOFR move independently, so the
dollar switch changed members' risk, which basis swaps restored.
\textbf{13.} Lost: its liability grew as discount factors rose.
\textbf{14.} The gain is roughly value times duration times 8.5 basis points,
and both the value of an in-the-money swap and its duration grow with
maturity; the thirty-year gain is 4.1 times the fifteen-year.
\textbf{15.} No: the compensation returns it to where it was. Its accounts may
show a gain and an equal cash payment, and a slightly different discount risk.
\textbf{16.} Mixed discounting within one netting pool would break the
equivalence between members' positions and the clearing house's margin, and a
single date lets compensation net across all accounts.
\textbf{17.} A bilateral amendment of the CSA to €STR (plus or minus a spread),
negotiated trade by trade, with its own compensation payment.
\textbf{18.} Little: EONIA was mechanically €STR plus 8.5 basis points, so the
switch was a known deterministic transfer, compensated in cash; there was no
unknown to trade.
\textbf{19.} Named result: \emph{the switch weekend} raised the value of the
fund's swap by EUR~1\,197\,426, which it paid back as cash compensation; after
the switch its discount sensitivity is EUR~141\,575 per basis point of €STR.
\textbf{20.} The rate paid on the collateral that secures it.
\end{solution}

\begin{solution}{iq:rc:multi-curve-and-collateral-discounting:1}
Because the trade is funded by its collateral: the holder of a positive-value
trade holds cash collateral equal to its value and pays the overnight rate on
it, so the value accrues at that rate. The index only determines the cash flows,
which are projected on its own curve.
\iqlookfor{discounting tied to the \omterm{def:rc:multi-curve-and-collateral-discounting:csa}{collateral rate}, projection to the index.}
\end{solution}

\begin{solution}{iq:rc:multi-curve-and-collateral-discounting:2}
First the overnight-index (€STR) \omterm{def:rc:multi-curve-and-collateral-discounting:curves}{discount curve} from OIS; then, holding it
fixed, the six-month Euribor \omterm{def:rc:multi-curve-and-collateral-discounting:curves}{projection curve} from the six-month fixing, the
forward rate agreements or futures, and swaps against six-month Euribor, each
priced with Euribor forwards discounted on €STR. The \omterm{def:rc:multi-curve-and-collateral-discounting:curves}{discount curve} comes first
because every projection instrument's price depends on it, not the reverse.
\iqlookfor{the order and the reason.}
\end{solution}

\begin{solution}{iq:rc:multi-curve-and-collateral-discounting:3}
Two floating legs of different tenors (six-month Euribor against €STR or
three-month Euribor) with a spread on one to make it fair. It is positive
because a six-month unsecured loan carries more credit and liquidity premium
than a rolled sequence of shorter or secured loans.
\iqlookfor{credit and liquidity premium by tenor.}
\end{solution}

\begin{solution}{iq:rc:multi-curve-and-collateral-discounting:4}
The poster will post whichever collateral is best for it at each date, the
one earning the highest rate once converted to the trade currency, so trades
are discounted at the pointwise maximum of the \omterm{def:rc:multi-curve-and-collateral-discounting:csa}{collateral rates}: receivables
are worth less and payables less negative. With volatile rates and basis the
option also has time value; it is often valued by simulation.
\iqlookfor{the maximum rule and the sign of the effect.}
\end{solution}

\begin{solution}{iq:rc:multi-curve-and-collateral-discounting:5}
Values change by the difference of the present values of the cash flows under
the two \omterm{def:rc:multi-curve-and-collateral-discounting:curves}{discount curves} (in-the-money receivers of fixed gain when rates fall);
discount risk changes from one index to the other. Fair compensation: cash
equal to the value change, computed with forwards held constant, netted per
account; and, where the two rates are not tied, basis swaps that restore the
old risk profile, with an auction for members who do not want them.
\iqlookfor{value and risk compensated separately.}
\end{solution}

\begin{solution}{iq:rc:multi-curve-and-collateral-discounting:6}
Par swaps on the \omterm{def:rc:multi-curve-and-collateral-discounting:curves}{projection curve}'s own index still reprice (both curves are
calibrated to them), so the error hides on off-market and seasoned trades,
where the discount rate is too high by the basis: an error of roughly value
times duration times the basis, 15 basis points here, EUR~60\,300 on a ten-year
swap 80 basis points in the money on EUR~100 million. A test with an off-market
trade finds it.
\iqlookfor{par trades hide the bug; off-market trades expose it.}
\end{solution}
